
\section{What this supplies to the construction program}
\label{nm15-ledger}

The new result supplies an actual nonlinear classical background
with an exact ordinary coarse constraint. The reduced stationary
equation includes both the nonlinear Wilson force and the full
curvature of the constraint manifold. It has uniform analytic
solutions, and the real action comparison survives removal of
the fixed gauge and transfer to ordinary averages. In this
homogeneous small-global-field class, the minimizer is no longer
merely an input that must be assumed.

Notebook 16 now derives a uniform local response and a controlled
expansion of the coarse classical action from these inputs.
Its weighted argument includes the derivative of the transposed
constraint map and the coarse multiplier. It proves factorial
responses with one weighted direction at every order and controls
the particular ordinary representative through its explicit gauge
formula. A fixed-exponent connected-tree bound for multiple
insertions is stronger and remains unproved.

Even a local classical coarse action would still be only part
of the required quantum renormalization step. Integration of
fluctuations brings determinants, gauge-coordinate Jacobians,
overlapping interactions and regions outside this global chart.
Those contributions must be bounded uniformly and iterated
with the coupling before a continuum probability law can be
claimed. None is replaced by the minimization theorem.

\begin{thebibliography}{9}
\bibitem{deposit15}
E.~Sussman and collaborators, \emph{Yang--Mills-LLM-experiment},
deposited raw accounts I--III at commit
\texttt{e392ef45013a581664bab76b90867170b30fe10e};
\href{https://github.com/ethanwsussman/Yang--Mills-LLM-experiment}{repository}.
Raw I:454--479 gives the Wilson normalization,
I:12756--12760 the physical bilinear pairing,
I:15471--15478 and I:15643--15648 the source quadratic
left-exponential chart, and I:43103--43107 the bare
constrained action. The ordinary constraint is displayed
at I:43098. Original journal papers are not freshly
certified by these deposited matches.
\bibitem{earlier15}
\emph{Toward a four-dimensional Yang--Mills construction},
research Notebooks 8--14 in the present workspace,
16 October 2026. Flat source-gauge extension, current and
regularity estimates; nonlinear constraint chart; nonflat
reference comparison and Wilson energy stability.
\end{thebibliography}
\end{document}
